% ===== Wrap-up (~4 min, then questions) =====
\partaccent{peach}\notopic

\begin{frame}{Այսօրվա 5 գաղափարը}
  \begin{enumerate}
    \item<1-> Զգայարանները տեսախցիկ չեն. ուղեղը «պղտոր առվակից» \textbf{վերակառուցում} է աշխարհը, ու սխալներն անխուսափելի են
    \item<2-> Հոտը ուղիղ կապված է հիշողության ու զգացմունքների հետ, իսկ համի մեծ մասը իրականում \textbf{հոտն} է
    \item<3-> Լսողությունն ու հպումը \textbf{մեխանիկական} են ու նույնիսկ գեներ են կիսում
    \item<4-> Մանրամասն տեսնում ենք միայն փոքրիկ \textbf{ֆովեայով}. մնացածը ուղեղը լրացնում ու գուշակում է, ու դեմքեր է տեսնում տոստի վրա
    \item<5-> Ուշադրությունը \textbf{սահմանափակ} է. hands-free-ն էլ է շեղում, ու կարող ենք չնկատել նույնիսկ, որ զրուցակիցը փոխվել է
  \end{enumerate}
\end{frame}

\begin{frame}{Փորձիր տանը}
  \begin{columns}[T]
    \begin{column}{0.31\textwidth}
      \begin{card}[colback=mint!30!cardwhite]
        \centering\bigemoji[1cm]{pencil}\ \bigemoji[1cm]{pencil}\par\smallskip
        \raggedright Երկու գրիչով, 3-4 սմ հեռավորությամբ, դիպչիր ընկերոջդ \textbf{մեջքին}։ Քանի՞ կետ կզգա։
      \end{card}
    \end{column}
    \begin{column}{0.31\textwidth}
      \begin{card}[colback=lavender!35!cardwhite]
        \centering\bigemoji[1cm]{phone}\par\smallskip
        \raggedright Կարդա մի տեքստ՝ \textbf{գլուխդ} թափահարելով։ Հետո գլուխդ պահիր անշարժ ու թափահարիր \textbf{հեռախոսը}։ Որ դեպքում է ավելի հեշտ։
      \end{card}
    \end{column}
    \begin{column}{0.31\textwidth}
      \begin{card}[colback=peach!35!cardwhite]
        \centering\bigemoji[1cm]{camera}\par\smallskip
        \raggedright Մի շաբաթ նկարիր \textbf{դեմքերը}, որ կտեսնես վարդակներում, մեքենաներում ու ուտելիքում։
      \end{card}
    \end{column}
  \end{columns}
\end{frame}

\begin{frame}{Հաջորդ վեբինարը}
  \begin{columns}[c]
    \begin{column}{0.6\textwidth}
      \chip[brainpink]{masks}{6. Անհատականություն}\par\medskip
      \begin{itemize}
        \item Երկաթե ձողը անցավ Ֆինեաս Գեյջի գլխով. իսկապե՞ս փոխվեց նրա բնավորությունը
        \item Որքանո՞վ կարելի է վստահել անհատականության թեստերին
        \item Ինչու է զայրույթը երբեմն օգտակար, ինչն է մեզ մոտիվացնում ու ինչու ենք ծիծաղում
      \end{itemize}
    \end{column}
    \begin{column}{0.34\textwidth}
      \centering\mascot[2.8cm]{angry}
    \end{column}
  \end{columns}
\end{frame}

\begin{frame}[plain]
  \begin{tikzpicture}[remember picture, overlay]
    \fill[lavender!60!cream] (current page.south west) rectangle (current page.north east);
    \fill[pattern={Dots[distance=4.4mm, radius=0.28mm]}, pattern color=plum!10!lavender!60!cream]
      (current page.south west) rectangle (current page.north east);
    \node at ([xshift=-3.2cm]current page.center) {\mascot[4.4cm]{happy}};
    \node[anchor=west, text=plum] at ([xshift=0.2cm, yshift=0.6cm]current page.center)
      {\titlefont\fontsize{40}{42}\selectfont Հարցեր};
    \node[anchor=west, text=plum, font=\mediumfont\large, align=left] at ([xshift=0.2cm, yshift=-1.0cm]current page.center)
      {Գրեք չաթում կամ\\ միացրեք խոսափողը \emoji{megaphone}};
    \node[anchor=west, text=plum!70!cream, font=\small] at ([xshift=0.2cm, yshift=-2.1cm]current page.center)
      {Դին Բըրնեթ, «The Idiot Brain», 2016};
  \end{tikzpicture}
\end{frame}
